\subsection{Random walk, and motion of a particle on a graph}
- Definire il random walk come processo stocastico discreto sul grafo.
- Definire l'equazione di master --> governa l'evoluzione della probabilità di trovare la particella in un certo nodo al tempo t.
- Far vedere che (in grafi non orientati e connessi) la probabilità stazionaria è proporzionale al grado del nodo.
\subsection{Transition matrix, jump probability}
- Deifnisco la matrice di transizione
- Mostrare la distribuzione di probabilità al tempo t
- Fare accenno che per tempi molto lunghi, il random walk è uequivalente all'equazione di diffusione (calore) governata dall'operatore laplaciano
\subsection{First passage time and mean recurrence time}
- First passage time: tempo richiesto per una particella per raggiungere un nodo target per la prima volta --> misura di una distanza "dinamica" tra nodi
- Mean recurrence time: tempo medio necessario per tornare la punto di partenza. Citare teorema di Kac  --> più un nodo è centrale , più velocemente ci si ritorna
- Mixing time --> tempo necessario per raggiungere la distribuzione stazionaria  --> per capire la velocità di reazione del sistema nervoso nel C. elegans
- Applicazioni --> usare questi trumenti per misurare l'efficenza della rete nel trasmettere segnali
\subsection{Case study: Signal propagation on the C. elegans connectome}
- Simulazione di un imissione di un segnale (random walk) in un neurone sensoriale
- Mostrare come il segnale di diffonde, il segnale rimarrà intrappolato nei cluster funzionali, prima di passare altrove?

Mi aspetto di ottenere --> neuroni hub (alta betweeneess) agiscono come stazioni di smistamento --> minimizzando il tempo di percorrenza globale.
Connectoma ottimizzato per avere un FPT molto basso tra input sensoriale e output motorio --> small world ??
- Fare confronto del tempo di percorrenza con un grafo casuale (Erdos-Renyi) con lo stesso numero di nodi. Usare per il confronto anche un null Model (grafo che mantenga la stessa distribuzione dei gradi --> per capire se le proprietà del verme dipendono dai neuroni hubb o dalla struttura specifica dei cluster)